\documentclass[8pt]{beamer}
\usepackage{tikz}
\usepackage{tikz-cd}
\usepackage{qtree}
\usepackage{amsfonts}
\usepackage{amssymb}
\usepackage{amsmath}
\usepackage{amscd}
\usepackage{amsthm}
\usepackage{tabularx}
% \usepackage{mathpartir}
\usepackage{bbm}
\usepackage[mathcal]{euscript}
% \usepackage[utf8]{inputenc}
\setbeamertemplate{navigation symbols}{}%remove navigation symbols
%\usepackage[T1]{fontenc}
\usepackage{lmodern} \title{Quantum theory of infinite systems}
\subtitle{from fields to information}
\author{Hans Halvorson}
\date[November 2013]{November 12, 2013} \usetheme{Copenhagen}
\begin{document}
\def\2{\mathcal}
\def\7{\mathbb}
\newenvironment{proenv}{\only{\setbeamercolor{local structure}{fg=green}}}{}
\newenvironment{conenv}{\only{\setbeamercolor{local structure}{fg=red}}}{}  
\newenvironment{neuenv}{\only{\setbeamercolor{local
      structure}{fg=yellow}}}{}

\newcommand{\3}{\mathfrak}

\begin{frame}[plain]
  \titlepage
\end{frame}

\begin{frame}{Elementary quantum mechanics}

\begin{description}

\item[states:] normalized vectors in a Hilbert space $\2H$.
  More generally: trace class operators on $\2H$.

\item[observables:] hermitian (self-adjoint) linear operators
  on $\2H$.

\item[statistical algorithm:] the expectation value of $O$ in state
  $\rho$ is given by 
$$\mathrm{Tr}(\rho O) \: = \: \sum _{i=1}^n\lambda _i\langle \psi _i,O\psi _i\rangle .$$

\item[dynamics] the Hamiltonian $H$ gives rise to a one-parameter
  group $U_t=\exp (-itH)$ of unitary operators on $\2H$.

\end{description}
\end{frame}

\begin{frame}{Elementary quantum mechanics}

  \begin{itemize} 

  \item The linear operators $\2B (\2H)$ on $\2H$ form an
    \textbf{algebra} closed under matrix sum, matrix multiplication,
    scalar multiplication, and the adjoint (hermitian conjugation)
    operation.
  \item This algebra has a natural \textbf{norm}: $\| O\|$ is the
    supremum of the eigenvalues of $O$.
  \item The states (density operators) form a \textbf{convex set}.
    The pure states correspond to unit vectors in $\2H$.
  \item The states (density operators) correspond to positive, linear
    maps from $\2B (\2H)$ to~$\7C$.
\end{itemize}

\end{frame}




\begin{frame}{$C^*$-algebras}

  \textbf{Definition:} An abstract $C^*$-algebra $\2A$ is a set
  equipped with operations:
\begin{itemize}
\item sum
\item product
\item complex scalar multiplication
\item adjunction
\item a norm $O\mapsto \| O\|$.
\end{itemize}


\end{frame}
\begin{frame}{$C^*$-algebras}

  \textbf{Fact:} Every finite-dimensional $C^*$-algebra is a direct
  sum of matrix algebras.

  But which $C^*$-algebras are not matrix algebras?  And why might we
  need them for quantum theory?

\end{frame}



\begin{frame}{states and probabilities}

  A \textbf{state} on $\2A$ is a normalized, positive, linear map from
  $\2A$ to $\7C$.

\bigskip The \textbf{transition probability} is defined by
$$ p(\omega ,\rho ) \: =\: 1-\frac{1}{4}\| \omega -\rho \| ^2 ,$$
where 
$$ \| \omega - \rho \| \: = \: \mathrm{sup}\{ |\omega (A)-\rho (A)| : A\in \2A, \| A\|\leq 1 \} .$$


\end{frame}

\begin{frame}{Example: infinite dimensional Hilbert spaces}

  Let $\2H$ be the linear space of square-integrable functions from
  $\7R ^n$ to $\7C$.  Then $\2H$ is an infinite dimensional Hilbert
  space.

  \bigskip A linear operator $O:\2H\to\2H$ is said to be
  \emph{bounded} if there is a $r<\infty$ such that $\| O\psi \|<r$
  for all unit-length $\psi$ in $\2H$.  The bounded linear operators
  $\2B (\2H )$ on $\2H$ form a $C^*$-algebra.


\end{frame}

\begin{frame}{Example: infinite spin chain}

  Objective: Form an infinite tensor product of the matrix algebra $M$

  \bigskip \textbf{observables:} the algebra $\2A = \bigotimes
  _{i\in \7Z} M_i$ is generated by operators
$$ \cdots \otimes A_{-2}\otimes A_{-1}\otimes A_0\otimes A_1\otimes A_2 \otimes \cdots $$
where $A_i=I$ for all but a finite number of entries.

\bigskip \textbf{states:} every vector $\otimes _{i\in\7Z} \psi _i$
induces a state on $\2A$.

\bigskip But there are too many states!  The Hilbert space would have
uncountably infinite dimension.  And there are too few observables!

\end{frame}

\begin{frame}{Representations of $C^*$-algebras}

  A $*$-\textbf{homomorphism} is a mapping $\pi :\2A\to\2B$ that
  preserves sum, product, and adjunction.

  \bigskip A \textbf{representation} of $\2A$ on $\2H$ is a
  $*$-homomorphism of $\2A$ into $\2B (\2H )$.

  \bigskip A $*$-homomorphism $\pi :\2A\to \2B$ gives rise to a mapping $\pi ^*:S(\2B )\to S(\2A )$.

$$ \2A \longrightarrow \2B \longrightarrow \7C $$

% \bigskip If $\pi$ is a representation of $\2A$ on $\2H$, then each
% density operator on $\2H$ induces a state of~$\2A$.



\end{frame}
\begin{frame}{Representations of $C^*$-algebras}

  \bigskip Two representations $(\2H ,\pi )$ and $(\2K ,\phi )$ of
  $\2A$ are said to be \textbf{unitarily equivalent} just in case
  there is a unitary mapping $U:\2H\to\2K$ such that
$$ U\pi (A) \: = \: \phi (A)U ,\qquad \forall A\in \2A .$$





\end{frame}


\begin{frame}{The GNS representation theorem}

  For each state $\omega$ of $\2A$, there is a representation $(\2H
  ,\pi )$ of $\2A$ and a unit vector $\Omega\in\2H$ such that
$$ \omega (A) \: =\: \langle \Omega ,\pi (A)\Omega \rangle ,\qquad \qquad \forall A\in \2A ,$$
and the vectors $\{ A\Omega : A\in \2A \}$ are dense in $\2H$.

\end{frame}


\begin{frame}{Von Neumann algebras}

A \textbf{von Neumann algebra} $\2M$ is a $*$-closed algebra acting on a Hilbert space $\2H$ and such that $\2M = \2M ''$, where 
$$ \2M ' \: = \: \{ X \in \2B (\2H ) : [X,A]=0, \forall A\in \2M \} .$$

\bigskip \textbf{Fact:} Every von Neumann algebra is a $C^*$-algebra.

\bigskip \textbf{Fact:} If $\2A$ is a $C^*$-algebra acting concretely
on $\2H$, then $\2A ''$ is the smallest von Neumann algebra containing
$\2A$.

\bigskip \textbf{Example:} Every matrix algebra is a von Neumann
algebra.  More generally, $\2B (\2H )$ is a von Neumann algebra.
\end{frame}

\begin{frame}{classification of von Neumann algebras}

\begin{tabularx}{\textwidth}{XXXX}
& type I & type II & type III \\ \\ 
\hline \\ 
application & elementary QM & QSM  & QFT
\end{tabularx}

\end{frame}


\begin{frame}{dynamics and operations}

  Each unitary operator $U\in \2A$ induces an inner automorphism
  $\mathrm{ad}_U:\2A\to\2A$ defined by
$$ \mathrm{ad}_U(A) = U^*AU,\qquad \forall A\in \2A .$$

\bigskip A \textbf{covariant representation} of $(\2A ,t\mapsto \alpha
_t)$ consists of a Hilbert space $\2H$, a $*$-homomorphism $\pi:\2A\to\2B (\2H )$, and a continuous representation $t\mapsto U_t$ of $\7R$ on $\2H$ such that 
$$ \pi (\alpha _t(A)) \: =\: U^*_t\pi (A)U_t ,\qquad \forall A\in \2A .$$

\bigskip \textbf{Dynamical GNS theorem:} If $\omega$ is an invariant
state of $(\2A ,t\mapsto \alpha _t)$, then the GNS representation is
covariant, and the GNS vector $\Omega$ is invariant under $U_t$.

\end{frame}
\begin{frame}{operations}

  A linear map $T:\2A\to\2B$ is said to be \textbf{positive} just in
  case $T(A)\geq 0$ for all $A\geq 0$, and normalized if $T(I)=I$.

  \bigskip A linear map $T:\2A\to\2B$ is said to be \textbf{completely
    positive} just in case for each $n\in\7N$ the map
  $T\otimes I_n:\2A\otimes M_n(\7C )\to\2A\otimes M_n(\7C )$ is
  positive.

\bigskip A map $T:\2B (\2H )\to \2B (\2H )$ is CP just in case there is a Hilbert space $\2K$ and a unitary map $U:\2H\otimes \2K\to \2H\otimes\2K$ such that 
$$ T(A) \: = \: \mathrm{Tr}(U^*(A\otimes I)U) .$$

\end{frame}

\begin{frame}{the SNAG theorem}

  Let $\mathbf{x}\to U(\mathbf{x})$ be a continuous representation of
  the translation group of $\7R ^{n+1}$ in the unitary operators on
  $\2H$.  Then there are $n+1$ self-adjoint operators $P_0,P_1,\dots
  ,P_{n}$ on $\2H$ such that
$$ U(\mathbf{x}) \: = \: \mathrm{exp}(-ix_0P_0)\mathrm{exp}(-ix_1P_1)\cdots \mathrm{exp}(-ix_{n}P_{n}) .$$

\bigskip The quadruple $(P_0,P_1,P_2,P_3)$ is called the 4-momentum of
the system.  The operator $P_0$ generates time translation, and is
called the Hamiltonian.

\end{frame}

\end{document}